\chapter{Aprendizaje por refuerzo mediante rúbricas extraídas y técnica de scaffolding}
\label{ch:metodo-propuesto}

El método propuesto busca adaptar modelos de lenguaje a dominios abiertos mediante una etapa de post training que combina adaptación supervisada y optimización por refuerzo con rúbricas extraídas de respuestas de referencia. La propuesta parte de que esas respuestas contienen información útil para construir una recompensa, aunque su calidad no pueda comprobarse íntegramente mediante una única regla determinista. En la implementación inicial, se extraen conceptos y expresiones relevantes, se organizan en criterios evaluables y se combinan procedimientos de verificación por código con una evaluación semántica para los casos no resueltos por coincidencia. A su vez, parte de esa información se presenta inicialmente como ayuda y se retira durante el entrenamiento, con el propósito de que el modelo aprenda a responder sin depender de ella.

El procedimiento se estudia mediante una evaluación inicial y un refinamiento posterior apoyado en datos de RaR. Las secciones siguientes distinguen los componentes comunes de las formulaciones propias de cada etapa y reúnen las arquitecturas y la configuración de entrenamiento en la sección~\ref{sec:metodo-entrenamiento}. El capítulo~\ref{ch:evolucion-entrenamiento} presenta los datos y el seguimiento del aprendizaje, mientras que el capítulo~\ref{ch:resultados-adaptacion} reúne la evaluación externa.

\section{Diseño general y separación de los datos}
\label{sec:metodo-datos}
El procedimiento parte de un conjunto de consultas y respuestas de referencia del dominio que se desea abordar. Antes del entrenamiento, los registros se reúnen en un formato común, se eliminan consultas duplicadas y se conserva su procedencia para poder vincular cada rúbrica con la respuesta de la que se extrajo.

En el esquema con una fuente común, el conjunto destinado al aprendizaje se divide en dos particiones: una destinada a la adaptación supervisada y otra reservada para construir las rúbricas y entrenar mediante GRPO. Así, si $\mathcal{D}$ representa el conjunto preparado, la separación se expresa como
\begin{equation}
\mathcal{D}=\mathcal{D}_{\mathrm{SFT}}\cup\mathcal{D}_{\mathrm{reserva}},
\qquad
\mathcal{D}_{\mathrm{SFT}}\cap\mathcal{D}_{\mathrm{reserva}}=\varnothing.
\label{eq:metodo-separacion}
\end{equation}
La primera partición permite adaptar el modelo siguiendo el procedimiento de la sección~\ref{sec:marco-sft}, mientras que las respuestas de la reserva aportan los conceptos utilizados para construir la señal de recompensa. Esta separación evita utilizar las mismas consultas en ambas etapas, aunque no garantiza por sí sola la ausencia de preguntas semánticamente similares.

Cuando se parte de un checkpoint supervisado existente, las fuentes de SFT y de GRPO pueden ser distintas, como ocurre con SFT45k y los datos procedentes de RaR. En ese caso, su independencia debe comprobarse sobre los registros disponibles, sin deducirla únicamente del nombre de las fuentes. Las proporciones de reparto son decisiones experimentales y no una condición fija de la metodología.

Por último, se reservan conjuntos de validación y prueba fuera del aprendizaje, mientras que la evaluación externa utiliza consultas y criterios independientes de los empleados en el entrenamiento, de modo que permite examinar las respuestas sin las ayudas proporcionadas durante GRPO. Las fuentes y las proporciones del particionado se detallan en la sección~\ref{sec:datos-experimentales}, mientras que las arquitecturas y el entrenamiento se reúnen en la sección~\ref{sec:metodo-entrenamiento}.
%
\section{Construcción y refinamiento de rúbricas a partir de referencias}
\label{sec:metodo-criterios}
En la primera etapa, se denominan entidades clave a los conceptos y expresiones relevantes de las respuestas de referencia que se utilizan como unidades de supervisión de la recompensa. El término incluye diagnósticos, localizaciones anatómicas, medicamentos, pruebas y profesionales, además de expresiones que describen acciones recomendadas o la evolución de un problema, de modo que no se restringe a entidades nombradas en sentido estricto. Los siguientes ejemplos presentan una selección real de expresiones de una de las rúbricas del conjunto central, conservando el idioma de los datos originales:

\begingroup
\begin{itemize}
\item \textbf{Boil:} representa un posible diagnóstico mencionado en la respuesta de referencia.
\item \textbf{pubic or groin area:} identifica la localización anatómica del problema descrito.
\item \textbf{resolved on its own:} describe la evolución del problema y su resolución espontánea.
\item \textbf{physical check:} recoge la evaluación física recomendada en la respuesta.
\item \textbf{Surgeon:} identifica al cirujano como uno de los profesionales mencionados.
\item \textbf{GP doctor:} identifica al médico general como otro de los profesionales mencionados.
\end{itemize}
\endgroup

Cada entidad se conserva con un identificador, un peso y una representación léxica que incluye formas normalizadas, lemas, alias y términos considerados esenciales para reconocerla. A partir de esa información se preparan criterios de presencia y, cuando corresponde, criterios de cantidades, postura, relaciones y acciones. Los primeros comprueban valores y unidades asociados a un concepto, mientras que la postura busca distinguir, por ejemplo, una posibilidad de una recomendación. Por su parte, las relaciones se aproximan mediante la presencia de conceptos en una misma oración y su cercanía en tokens, y las acciones se reconocen mediante expresiones compatibles con la familia de acción esperada.

Estas comprobaciones buscan conservar algo más que una lista de palabras, aunque su implementación sigue siendo una aproximación. La aparición de una expresión de posibilidad en una oración no garantiza que modifique el diagnóstico correcto, del mismo modo que dos conceptos cercanos no necesariamente mantienen la relación clínica esperada. Además, la extracción puede conservar fragmentos poco pertinentes, como ocurre en la misma muestra con una referencia a una consulta de seguimiento de pago. Por ello, la selección de entidades también debe revisarse y no puede asumirse que toda expresión extraída constituye un criterio clínico útil.


En la segunda etapa, la construcción parte de las obligaciones explícitas de la pregunta y utiliza la respuesta de referencia para fundamentar los criterios. El objetivo es exigir el contenido necesario para resolver la consulta, sin convertir cada detalle de la referencia en un requisito adicional. Por ello, se identifica si la tarea solicita una respuesta directa, una explicación o varias acciones y se vincula cada criterio con la obligación que evalúa.

Cada criterio evalúa un requisito concreto de la consulta y se vincula con un fragmento de la referencia. También indica qué respuestas equivalentes se aceptan y qué menciones resultarían insuficientes. Por ejemplo, si la consulta pide identificar una sustancia cuyo nivel está elevado, nombrarla sin indicar el sentido del cambio no cubre el requisito completo. El criterio conserva esa relación y su polaridad, y registra la certeza, las condiciones temporales o los parámetros cuando son necesarios para interpretar la respuesta. Si la referencia no permite fundamentar un requisito solicitado, se registra la falta de cobertura.

La generación y la revisión mediante modelos de lenguaje se acompañan de comprobaciones estructurales y de correspondencia con la evidencia. Este control busca mejorar la fidelidad de la rúbrica a la consulta y a la referencia, aunque no constituye una revisión médica de su corrección. Los criterios aceptados se clasifican como esenciales o adicionales y se puntúan con la formulación correspondiente a la segunda etapa, descrita en la sección~\ref{sec:metodo-recompensa}.
%
\section{Evaluación de criterios y formulaciones de recompensa}
\label{sec:metodo-recompensa}
\subsection{Evaluación híbrida y recompensa de la primera etapa}
La evaluación inicial comienza con un procedimiento determinista que normaliza la respuesta y busca las entidades dentro de sus oraciones. Una coincidencia puede obtenerse mediante una expresión exacta o mediante una cobertura suficiente de grupos de términos y anclas, utilizando los lemas y los alias previstos para cada concepto. Este procedimiento reduce el costo de evaluación y permite registrar qué términos respaldan cada detección, aunque puede producir falsos negativos ante paráfrasis y falsos positivos cuando una expresión aparece en un contexto inadecuado.

Para ampliar la cobertura, las entidades no detectadas por código se envían a un juez semántico junto con la consulta, la respuesta generada y la información asociada al criterio. El juez debe señalar si el concepto está presente y aportar un fragmento de evidencia, cuya presencia en una oración de la respuesta se comprueba antes de aceptar la detección. De este modo, la evaluación semántica puede recuperar menciones que el diccionario no reconoce, pero no reemplaza todas las comprobaciones ni revisa automáticamente las coincidencias positivas del código. La verificación del fragmento acredita que la evidencia está en el texto, sin garantizar por sí sola la corrección clínica de la interpretación.

La recompensa combina cinco componentes, cuyos valores se denotan por $E$, $N$, $T$, $B$ y $A$, correspondientes a entidades, cantidades, postura, relaciones y acciones. Como algunos ejemplos no contienen todos los tipos de criterio, los pesos se normalizan entre los componentes disponibles. Si $\mathcal{K}_i$ reúne los componentes activos del ejemplo $i$ y $b_k$ es el presupuesto del componente $k$, se define
\begin{equation}
\alpha_{ik}=\frac{b_k}{\sum_{\ell\in\mathcal{K}_i}b_\ell},
\qquad
R_i(y)=\operatorname{clip}\!\left(
\sum_{k\in\mathcal{K}_i}\alpha_{ik}C_{ik}(y)-P_i(y),-1,1\right),
\label{eq:metodo-recompensa}
\end{equation}
donde $C_{ik}(y)$ representa el valor del componente evaluado sobre la respuesta $y$ y $P_i(y)$ reúne las penalizaciones. Así, la ausencia de criterios numéricos en un ejemplo no se interpreta como un incumplimiento ni reduce automáticamente su recompensa máxima.

Las penalizaciones contemplan respuestas vacías, repetición de secuencias y listas breves de palabras clave, buscando reducir formas de obtener cobertura sin desarrollar una explicación útil. Los presupuestos de los componentes y los valores de las penalizaciones utilizados en el experimento central se detallan en la sección~\ref{sec:protocolo-configuracion}. Estos mecanismos no constituyen una evaluación completa de seguridad clínica, ya que no comprueban de manera general contraindicaciones, daño potencial ni la pertinencia de toda recomendación.


\subsection{Evaluación semántica y recompensa del piloto posterior}
En el piloto de la segunda etapa, el juez evalúa el significado de cada criterio y asigna un estado de cumplimiento, ausencia, contradicción, no aplicabilidad o incertidumbre. La decisión debe estar respaldada por evidencia literal cuando corresponde, y una contradicción aplicable impide considerar satisfecho el criterio. Así, el control por código comprueba la integridad del veredicto y la evidencia, mientras que el juicio sobre el cumplimiento del requisito es semántico; este reparto difiere de la búsqueda léxica seguida de revisión de entidades no detectadas utilizada inicialmente.

Si $\mathcal{A}_i$ es el conjunto de criterios aplicables a la respuesta $y$ del ejemplo $i$, $v_j$ su peso positivo y $s_j(y)$ vale 1 únicamente cuando el criterio está satisfecho, la recompensa del piloto es
\begin{equation}
R_i^{\mathrm{piloto}}(y)=
\frac{\sum_{j\in\mathcal{A}_i}v_j s_j(y)}{\sum_{j\in\mathcal{A}_i}v_j}.
\label{eq:metodo-recompensa-piloto}
\end{equation}
Los criterios esenciales tienen peso 3 y los adicionales peso 1. La ausencia o la contradicción no aportan crédito y no reciben una penalización numérica adicional, mientras que la no aplicabilidad requiere una condición prevista en la rúbrica. Por su parte, la incertidumbre, la evidencia inválida o los veredictos incompletos bloquean el cálculo y la actualización correspondiente, en lugar de convertirse automáticamente en una recompensa nula.

Esta formulación concreta permite estudiar una señal basada en la satisfacción de obligaciones y debe distinguirse de otras variantes de recompensa que se investiguen posteriormente. El contraste con RaR requiere examinar tanto la construcción de los criterios como su agregación~\cite{gunjal2025rar}; utilizar preguntas y respuestas de ese trabajo no implica conservar su rúbrica, sus pesos ni su protocolo experimental.
%
\section{Retirada progresiva de ayudas durante el entrenamiento}
\label{sec:metodo-ayudas}
La retirada progresiva de ayudas, denominada scaffold, presenta al modelo parte de las entidades esperadas junto con la consulta al comienzo del entrenamiento. Su propósito es facilitar la incorporación del contenido del dominio y aumentar después la exigencia de recuperación autónoma, por lo que las entidades visibles se reducen a medida que avanza el calendario. La rúbrica completa permanece disponible para evaluar la respuesta, aunque el modelo recibe únicamente la fracción de conceptos habilitada para esa presentación.

El calendario establece una fracción de entidades visibles para cada etapa, que se reduce hasta llegar a consultas sin ayudas. Para ello, se asigna un orden a las entidades de cada ejemplo y se conservan subconjuntos cada vez menores, manteniendo la rúbrica completa para evaluar todas las respuestas. La duración del calendario y las reglas de selección empleadas en el experimento central se detallan en la sección~\ref{sec:protocolo-configuracion}.

La figura~\ref{fig:scaffolding-etapas} ilustra este mecanismo con el calendario definido en el diseño posterior V29. En esa implementación, un modelo de lenguaje genera ayudas breves a partir de los criterios y de la evidencia de la respuesta de referencia, y una revisión mediante LLM comprueba su correspondencia con esas fuentes antes de incorporarlas al entrenamiento. Las ayudas aceptadas se ordenan de manera determinista y, en cada etapa, se conserva la cantidad correspondiente a la tasa de retención, redondeada hacia arriba. Por ejemplo, cuatro ayudas dan lugar a subconjuntos de cuatro, tres, dos, una y ninguna ayuda para las tasas de 100\%, 75\%, 50\%, 25\% y 0\%, respectivamente. Este ejemplo representa el diseño de V29; el calendario del experimento central evaluado en este informe se especifica en la sección~\ref{sec:protocolo-configuracion}.

\begin{figure}[htbp]
\centering
\begingroup
\singlespacing
\shorthandoff{<>}
\definecolor{scaffoldink}{HTML}{253F53}
\definecolor{scaffoldline}{HTML}{658AA8}
\definecolor{scaffoldcream}{HTML}{F7FAFC}
\definecolor{scaffoldpeach}{HTML}{C7DCEC}
\definecolor{scaffoldsand}{HTML}{DBE7F0}
\definecolor{scaffoldshadow}{HTML}{D6E0E8}
\resizebox{\textwidth}{!}{%
\begin{tikzpicture}[>=Stealth,
 base/.style={rounded corners=9pt,draw=scaffoldline,line width=.9pt,fill=scaffoldcream},
 title/.style={align=center,text=scaffoldink,font=\sffamily\small\bfseries,inner sep=2pt},
 body/.style={align=center,text=scaffoldink,font=\sffamily\small,inner sep=0pt},
 branch/.style={draw=scaffoldline,line width=1.3pt,rounded corners=8pt,line cap=round,line join=round}]
% Tarjetas superiores: cabecera y cuerpo, con sombra discreta.
\fill[scaffoldshadow,rounded corners=9pt] (4.04,4.23) rectangle (9.24,5.71);
\draw[base] (4.12,4.32) rectangle (9.32,5.8);
\draw[draw=scaffoldline,fill=scaffoldpeach,rounded corners=6pt,line width=.8pt]
 (4.3,5.15) rectangle (9.14,5.59);
\node[title] at (6.72,5.37) {LLM};
\node[body] at (6.72,4.73) {Generación de rúbricas};
\fill[scaffoldshadow,rounded corners=9pt] (4.04,2.48) rectangle (9.24,3.61);
\draw[base] (4.12,2.57) rectangle (9.32,3.7);
\draw[draw=scaffoldline,fill=scaffoldsand,rounded corners=6pt,line width=.8pt]
 (4.3,2.88) rectangle (9.14,3.39);
\node[title] at (6.72,3.135) {Entrenamiento con GRPO};
\draw[branch] (6.72,4.32) -- (6.72,3.7);
% Una única línea conecta GRPO con todas las etapas.
\draw[branch] (6.72,2.57) -- (6.72,1.92);
\draw[branch] (0,1.12) -- (0,1.64) -- (13.44,1.64) -- (13.44,1.12);
\foreach \x in {3.36,6.72,10.08} {
 \draw[branch] (\x,1.64) -- (\x,1.12);
}
\draw[branch] (6.72,1.92) -- (6.72,1.64);
% Etapas con cabeceras de tonos azules y porcentajes destacados.
\foreach \x/\n/\pct/\tone in {0/1/100/scaffoldpeach,3.36/2/75/scaffoldpeach!75!scaffoldsand,6.72/3/50/scaffoldsand,10.08/4/25/scaffoldsand!65!scaffoldcream,13.44/5/0/scaffoldsand!35!scaffoldcream} {
 \fill[scaffoldshadow,rounded corners=9pt] (\x-1.5,-1.76) rectangle (\x+1.34,1.03);
 \draw[base] (\x-1.42,-1.67) rectangle (\x+1.42,1.12);
 \draw[draw=scaffoldline,fill=\tone,rounded corners=6pt,line width=.8pt]
 (\x-1.25,.46) rectangle (\x+1.25,.91);
 \node[title] at (\x,.685) {Etapa \n};
 \node[body,font=\normalfont\large] at (\x,-.13) {\pct\%};
 \node[body] at (\x,-1.02) {de entidades\\en el input};
}
\end{tikzpicture}%
}
\endgroup

\caption{Retirada progresiva de ayudas en cinco etapas de entrenamiento.}
\label{fig:scaffolding-etapas}
\end{figure}

En la formulación inicial, la recompensa distingue entre conceptos visibles y ocultos para dar mayor peso a la recuperación sin ayuda. Si $w_e$ es el peso base de la entidad $e$, $m_e$ su puntuación de detección y $\lambda_e$ es un multiplicador que toma un valor mayor para las entidades ocultas que para las visibles, el componente de entidades se calcula como
\begin{equation}
E_i(y)=\min\!\left(1,
\frac{\sum_e w_e\lambda_e m_e}{\sum_e w_e\lambda_e}
+\eta\,\frac{H_i^{\mathrm{rec}}}{H_i}\right),
\label{eq:metodo-entidades-ocultas}
\end{equation}
donde $H_i$ es la cantidad de entidades ocultas, $H_i^{\mathrm{rec}}$ la cantidad recuperada y $\eta$ controla la bonificación por su recuperación. El término adicional se toma como cero cuando no hay entidades ocultas. Esta formulación añade una bonificación máxima de $\eta$ al componente, que luego se combina con los restantes mediante la ecuación~\ref{eq:metodo-recompensa}.

El mecanismo busca que recuperar información no presentada resulte más valioso que repetir las ayudas, aunque su efecto depende de la composición de cada ejemplo. Cuando todas las entidades están ocultas, el multiplicador común se cancela en el cociente ponderado y la diferencia adicional procede de la bonificación. Por ello, la asimetría tiene un efecto relativo especialmente en las etapas con entidades visibles y ocultas. Asimismo, las respuestas generadas con ayudas deben distinguirse de las producidas sin ellas al analizar el comportamiento del modelo.


En el piloto posterior, las ayudas se vinculan a los criterios y a la evidencia de referencia y se retiran siguiendo el calendario de cinco etapas de la figura~\ref{fig:scaffolding-etapas}. Las tasas representan retención nominal: el redondeo hacia arriba puede producir la misma cantidad de ayudas en etapas consecutivas cuando hay pocas disponibles. La validación se realiza sin ayudas, mientras que el juez utiliza la rúbrica completa y no recibe las ayudas visibles en el prompt del modelo. A diferencia de la formulación inicial, la recompensa de la ecuación~\ref{eq:metodo-recompensa-piloto} conserva los pesos de los criterios durante las etapas y no aplica el multiplicador por entidades ocultas de la ecuación~\ref{eq:metodo-entidades-ocultas}.
%
\section{Integración de las rúbricas y las ayudas con GRPO}
\label{sec:metodo-integracion}
Una vez preparado el conjunto con sus rúbricas y máscaras de visibilidad, el modelo adaptado genera un grupo de respuestas para cada consulta y cada respuesta recibe una recompensa mediante el procedimiento anterior. Esos valores alimentan la comparación relativa descrita en la sección \ref{sec:grpo}, de modo que GRPO actualiza la política utilizando las diferencias dentro del grupo~\cite{shao2024deepseekmath}. La propuesta conserva ese mecanismo de optimización y modifica principalmente la construcción de la señal y las ayudas disponibles durante las distintas etapas.

Durante el entrenamiento se registran las respuestas, las decisiones de presencia, la evidencia semántica, los valores de cada componente y las penalizaciones. Estos registros permiten analizar si una variación de la recompensa procede de recuperar contenido, de cambios en las ayudas o de evitar determinadas formas de respuesta. La configuración utilizada para aplicar este procedimiento se describe en la sección~\ref{sec:protocolo-configuracion}, mientras que la evaluación externa emplea criterios independientes de los utilizados para entrenar.


La misma secuencia de generación, evaluación y actualización se mantiene en la segunda etapa, utilizando las rúbricas refinadas y la recompensa del piloto. En consecuencia, los cambios en la construcción de criterios y en su puntuación pueden estudiarse dentro del procedimiento común, conservando una referencia KL correspondiente al modelo adaptado mediante SFT. Las configuraciones de ambas etapas se detallan en la sección~\ref{sec:metodo-entrenamiento}.
%
\section{Entrenamiento}
\label{sec:metodo-entrenamiento}
\label{sec:metodo-evolucion}
Las pruebas iniciales utilizan Qwen2.5-7B, mientras que la comparación con el trabajo externo parte de Qwen2.5-3B SFT45k.

\subsection{Elección de las arquitecturas}
\label{sec:metodo-arquitecturas}
Se utilizan dos tamaños de la familia Qwen2.5, de 7B y 3B parámetros~\cite{qwen2024qwen25}. El alcance de las comparaciones se establece en la sección~\ref{sec:resultados-modelos}.

En la primera etapa, Qwen2.5-7B Base proporciona el punto de partida para obtener el modelo SFT-90 y estudiar después su adaptación con GRPO. Durante GRPO se actualizan los parámetros de LoRA asociados a las proyecciones de atención y del bloque de transformación, mientras que los embeddings y la cabeza de salida guardados con el adaptador se congelan en esta etapa.

Para estudiar las nuevas rúbricas y recompensas se utiliza Qwen2.5-3B, buscando reducir el costo de los experimentos. Se parte del checkpoint SFT45k adaptado mediante QLoRA y actualización de embeddings. El adaptador supervisado se fusiona con el modelo base y se conservan los embeddings adaptados; a partir de esa misma versión se construyen la política entrenable y la referencia congelada utilizada para la regularización KL. Durante GRPO se incorporan nuevos adaptadores LoRA y se entrenan también los embeddings, manteniendo el enlace con la cabeza de salida cuando corresponde.

\subsection{Configuración de entrenamiento}
\label{sec:metodo-configuracion}
\label{sec:protocolo-configuracion}
\subsubsection{Pruebas del método de adaptación a dominios abiertos}
En la primera etapa, la adaptación supervisada utiliza la partición del 90\% de las consultas y respuestas preparadas, mientras que la reserva aporta los ejemplos para las rúbricas y GRPO. El entrenamiento posterior parte de SFT-90 y actualiza adaptadores LoRA de rango 64 y escala 16. Las fuentes y el particionado se detallan en la sección~\ref{sec:protocolo-primera-etapa}.

El entrenamiento utiliza 16 generaciones por consulta, una tasa de aprendizaje de $5\times10^{-6}$, un coeficiente de regularización KL de 0,04 y un máximo de 250 tokens nuevos por respuesta, con temperatura 0,7 y top\_p de 0,9. Cada una de las 500 consultas se presenta en diez etapas de ayudas, formando un conjunto de 5.000 entradas que el entrenador recorre una vez.

Los pesos iniciales de los componentes de recompensa son 0,55 para entidades, 0,20 para cantidades, 0,15 para postura, 0,05 para relaciones y 0,05 para acciones, normalizados entre los componentes activos de cada ejemplo según la ecuación~\ref{eq:metodo-recompensa}. Las penalizaciones corresponden a respuestas vacías, repetición de secuencias de tres tokens y listas breves de palabras clave, con un valor de 1 para las respuestas vacías y máximos de 0,10 y 0,05 para las otras dos, respectivamente.

Durante las primeras nueve etapas, todas pueden recibir ayudas, pero la proporción de entidades conservadas como ayuda desciende del 100\% al 20\% en pasos de diez puntos porcentuales, mientras que en la décima etapa no se muestra ninguna entidad. La cantidad visible se obtiene redondeando hacia arriba y conservando al menos una entidad cuando existe asistencia, de modo que la proporción efectiva depende de cuántos conceptos contenga cada rúbrica. Además, la selección utiliza un orden determinista que permite que las ayudas de una etapa posterior formen un subconjunto de las anteriores.

En la ecuación~\ref{eq:metodo-entidades-ocultas}, el multiplicador $\lambda_e$ vale 1 para las entidades visibles y 2,5 para las ocultas, mientras que la bonificación máxima por recuperación se fija en $\eta=0{,}05$.

\subsubsection{Comparación con un trabajo externo}
El entrenamiento supervisado de SFT45k comprende dos épocas sobre 45.000 muestras. El entrenamiento con GRPO parte de este checkpoint y utiliza preguntas y respuestas procedentes de RaR, cuyos conjuntos de entrenamiento y validación se describen en la sección~\ref{sec:protocolo-segunda-etapa}.

El piloto configura nuevos adaptadores LoRA de rango 16 y escala 32, con una tasa de aprendizaje de $2\times10^{-6}$, mientras que los embeddings se actualizan con una tasa de $2\times10^{-7}$. Se generan 16 respuestas por consulta con temperatura 0,7, top\_p de 0,9 y un máximo de 1.024 tokens nuevos. El coeficiente de regularización KL es 0,04 y el objetivo utiliza normalización BNPO, con una única actualización por grupo. La generación se organiza en microbatches de ocho respuestas y el cálculo de gradientes en cuatro microbatches de cuatro respuestas.

El calendario del piloto recorre las consultas en cinco etapas con proporciones de ayudas de 100\%, 75\%, 50\%, 25\% y 0\%, utilizando el orden fijo y el redondeo hacia arriba descritos en la sección~\ref{sec:metodo-ayudas}. La validación se realiza sin ayudas antes del entrenamiento y al finalizar cada etapa. La puntuación sigue la ecuación~\ref{eq:metodo-recompensa-piloto}, con pesos de 3 para criterios esenciales y 1 para adicionales, y permanece independiente de las ayudas visibles para el modelo.

%
\section{Métricas de evaluación}
\label{sec:metodo-metricas}
La pérdida supervisada permite seguir el ajuste a las respuestas de referencia, mientras que la recompensa de GRPO mide el cumplimiento de las rúbricas de entrenamiento. La evaluación en HealthBench examina después las respuestas con criterios externos al aprendizaje.

\subsection{Seguimiento del ajuste supervisado}
\label{sec:metricas-sft}
Durante SFT se registra la pérdida de entrenamiento sobre los tokens supervisados, definida en la ecuación~\ref{eq:sft-mascara}. Su descenso indica que el modelo asigna mayor probabilidad a las respuestas de referencia. La pérdida de validación aplica el mismo cálculo a ejemplos separados, sin actualizar parámetros, para seguir el ajuste fuera de los datos utilizados en cada paso.

Los registros también incluyen la tasa de aprendizaje y la norma del gradiente. En SFT45k se contabilizan los tokens supervisados y su velocidad de procesamiento, lo que permite relacionar cada actualización con el volumen de texto aprendido y el tiempo empleado. Estos valores describen el comportamiento del entrenamiento; la calidad de las respuestas generadas se examina mediante la evaluación posterior.

\subsection{Seguimiento del entrenamiento con GRPO}
\label{sec:metricas-grpo}
La medida principal es la recompensa media de las respuestas generadas para cada consulta. Junto con ella se registra su desviación estándar dentro del grupo, porque las diferencias entre respuestas determinan las ventajas utilizadas por GRPO. En el piloto posterior también se identifica la frecuencia de grupos con recompensas idénticas. Para visualizar la tendencia se conserva una media móvil exponencial, además de los valores sin suavizar.

La pérdida de optimización, la divergencia KL respecto de la referencia y la norma del gradiente permiten seguir las actualizaciones. También se registra la longitud de las respuestas; en el piloto se añade la proporción de salidas que alcanzan el límite de generación. Las métricas se relacionan con la etapa y con la cantidad de ayudas visibles, ya que una recompensa obtenida con asistencia completa no tiene la misma dificultad que una obtenida sin ayudas.

En el piloto posterior se registra por separado la proporción de criterios esenciales y adicionales satisfechos, junto con los estados asignados por el juez. La validación utiliza consultas sin ayudas y registra la recompensa media, su dispersión y la proporción de respuestas truncadas, antes del entrenamiento y al terminar cada etapa. Esta evaluación permite observar el comportamiento sin asistencia mientras cambia la cantidad de ayudas del conjunto de entrenamiento.

\subsection{Evaluación externa en HealthBench}
\label{sec:metricas-benchmark}
HealthBench combina los criterios cumplidos y sus pesos en una puntuación por conversación. El resultado global promedia esos valores y aplica el recorte definido en la ecuación~\ref{eq:healthbench-global}~\cite{arora2025healthbench,openai2026healthbenchcode}. El procedimiento se desarrolla en la sección~\ref{sec:marco-evaluacion}, mientras que las condiciones del juez y de generación del proyecto se detallan en la sección~\ref{sec:protocolo-experimental}.

Para analizar la primera etapa se presentan desgloses de precisión, calidad de comunicación, comprensión del contexto y seguimiento de instrucciones, calculados según la ecuación~\ref{eq:healthbench-ejes}. También se cuentan los criterios positivos cumplidos y los negativos activados, con sus respectivas proporciones. Estos conteos permiten reconocer si una mayor cobertura se acompaña de más errores, aunque no incorporan los pesos utilizados en la puntuación del benchmark.

La comparación entre SFT-90 y SFT-90 + GRPO utiliza los mismos casos y expresa las diferencias entre ejes en puntos porcentuales. Los fragmentos de respuestas del análisis cualitativo ayudan a interpretar qué cambios de contenido acompañan esas diferencias. Los resultados corresponden a la evaluación con el juez local utilizado en el proyecto.


